\section*{Part III: Candidate Sourcing, Access, and the Creation of a Real Choice Set}

\subsection*{23. Selection Begins Before an Application Is Submitted}

An employer can choose only among people who know the opportunity exists, consider it credible, find it relevant, and are willing to apply. The set of applicants is therefore not a random sample of human talent. It is produced by the employer's brand, network, salary, location, terms, role design, channel choices, and how people perceive the organization. An excellent selection battery applied to an extremely narrow pool may identify the best person in that pool while missing most of the people who could have succeeded.

This matters especially for a new organization without a famous name. Familiar employers can attract applicants through recognition and perceived security. Unknown ventures must provide more information about the work, the founder's competence, compensation reliability, expected support, and purpose. A small institute should not confuse low application volume with proof that the occupation lacks suitable talent. It may have a discoverability or credibility problem.

Hiring therefore contains 2 separate optimization problems: \textbf{construct a sufficiently capable and diverse set of candidates} and \textbf{evaluate them using job-relevant evidence}. Improvement in 1 cannot fully compensate for failure in the other.

\subsection*{24. Sourcing Channels Produce Different Applicant Populations}

Employee referrals, alumni networks, professional communities, public job boards, university placements, recruiter searches, open-source contributors, specialized forums, and internal mobility expose the employer to different people. They also create different biases and costs. Personal referrals may improve access to credible background information but concentrate opportunity inside a familiar social network. Large job boards create volume but may require extensive screening. Specialist communities can reveal work publicly but may underrepresent capable people who do not participate in that community.

A founder should classify channels by 5 properties: access to role-relevant talent, applicant interest, candidate diversity, screening cost, and likelihood of a timely hire. Use a small channel portfolio rather than pretending there is 1 universally superior source. For a technical educator, professional teaching communities and technical-practice groups may offer different candidates; the former reveal instructional experience, while the latter reveal subject expertise.

The useful measurement is not merely \textit{applications per channel} but \textbf{qualified applicants and successful hires per unit of recruitment cost and applicant burden}, with longer-term evaluation once actual performance becomes observable.

\subsection*{25. Referrals Are Information Advantages and Potential Exclusion Mechanisms}

A thoughtful referral from someone who has closely observed relevant work can provide information unavailable from a résumé. ``I have watched her diagnose a failed deployment and coordinate an incident response'' is useful. ``He is my college friend and a great guy'' provides little job evidence. Social proximity can improve access, but its informational value depends on what the referrer actually knows.

Referral-heavy recruitment may also reproduce existing class, geographic, educational, and social networks. Even without discriminatory intent, the employer may repeatedly see candidates from the same circles. The founder should therefore ask referrers to describe observed work and should keep at least 1 open route through which unfamiliar applicants can be considered on comparable terms.

Referrals should usually accelerate discovery, not automatically bypass critical evaluation. An exception could be an urgent short-term engagement with tightly bounded risk and an unusually strong direct work history. The organization must still know the basis for trusting the recommendation and should document any material exemption from its ordinary process.

\subsection*{26. An Unknown Employer Must Sell the Work, Not Manufacture Prestige}

Strong candidates have alternatives. They assess the employer's likely stability, quality of leadership, autonomy, learning opportunities, work conditions, and whether promises will be honored. An unknown institution can attract unusually serious candidates by presenting a consequential, well-defined problem and visibly competent leadership. It cannot ethically imply benefits it has not earned: world-renowned faculty colleagues, guaranteed rapid promotion, a mature research environment, or venture-scale upside without a genuine agreement.

A credible early recruiting proposition might say: \textit{The institution is small; the systems are being built. The work involves developing rigorous project-review standards and seeing how they affect actual learner work. The initial engagement is part-time and paid. Teaching quality matters more than public profile. Here is the sample work and how your contribution will be assessed.}

That proposition will repel some applicants who want a stable large institution. This is useful information. The goal is not universal appeal but an accurate meeting between the role's conditions and the candidate's preferences and circumstances.

\subsection*{27. Sourcing Is Constrained by Compensation and Employment Design}

A startup offering an unpredictable 6-hour weekly contract cannot expect the same candidate pool as a firm offering a secure full-time salary. A role requiring evening and weekend availability may exclude otherwise excellent applicants with caregiving commitments. An international posting may create employment-law, tax, currency, and data-management complications. Equity can matter to candidates only when the terms, risk, and potential dilution are real and understood.

The founder should choose between employment, consulting, project-based contracting, adjunct engagement, and fractional roles based on \textbf{the actual nature of work and applicable law}, not merely on avoiding benefits or obligations. Classifying someone as a contractor does not automatically make them legally independent. Employment status, confidentiality, intellectual property, taxation, and payment obligations require appropriate local advice.

Selection criteria must match the engagement. A part-time reviewer may need strong asynchronous judgment and reliable turnaround but not a desire to manage a department. Asking them to demonstrate executive strategy may be irrelevant and costly. \textbf{Compensation, job definition, and screening should describe the same exchange.}
